\documentclass[12pt]{article}

% --- Packages ---
\usepackage[a4paper, margin=2.5cm]{geometry}
\usepackage{amsmath}
\usepackage{amssymb}
\usepackage{amsthm}
\usepackage{graphicx}
\usepackage[colorlinks=true, linkcolor=blue, citecolor=blue, urlcolor=blue]{hyperref}
\usepackage{booktabs}
\usepackage{natbib}

% --- Theorem Environments ---
\newtheorem{theorem}{Theorem}
\newtheorem{definition}[theorem]{Definition}

% --- Title ---
\title{The Crystallisation Operator: Mathematical Structure and Properties}
\author{%
  Lee Smart\\[4pt]
  \textit{Institute of Vibrational Field Dynamics}\\[2pt]
  \texttt{contact@vibrationalfielddynamics.org}\\[2pt]
  \texttt{@vfd\_org}%
}
\date{March 2026}

\begin{document}

\maketitle

\begin{abstract}
We introduce the \textbf{crystallisation operator}, a mathematical framework that replaces the discontinuous, probabilistic notion of wavefunction collapse with a deterministic constraint-driven process. In this formulation an unresolved multi-state configuration evolves under lawful constraints toward a stable, coherent structure by minimising a closure functional that balances constraint satisfaction, energetic cost, and phase coherence. The observed definite outcome is reinterpreted not as a random projection but as the formation of a closure-compatible eigenstate --- a \textit{crystal} in configuration space. We develop the operator definition, its dynamical flow form, a spectral reweighting formulation, and prove existence, stability, and mode-selection theorems. Applications to quantum measurement are discussed, with an outlook on broader applicability.
\end{abstract}

%------------------------------------------------------------------
\section{Introduction}
%------------------------------------------------------------------

The measurement problem in quantum mechanics is conventionally addressed through wavefunction collapse: upon measurement, a superposition is said to discontinuously reduce to a single eigenstate with probability given by the Born rule. Despite its empirical success, this prescription suffers from several conceptual difficulties:
\begin{itemize}
    \item It introduces a fundamental discontinuity absent from unitary evolution.
    \item The selection mechanism remains unspecified --- \textit{why} a particular outcome obtains is left to irreducible randomness.
    \item The boundary between ``quantum'' and ``classical'' regimes is vague.
\end{itemize}

The Vacuum Field Dynamics (VFD) programme offers an alternative perspective. Rather than treating definite outcomes as the result of stochastic projection, VFD frames state resolution as a \textbf{physical process of crystallisation}: the system evolves from an unstable, over-complete, constraint-incompatible configuration into a stable, constraint-satisfying coherent structure.

This paper formalises that intuition.

%------------------------------------------------------------------
\section{Collapse vs Crystallisation}
%------------------------------------------------------------------

\subsection{Standard Collapse}

In orthodox quantum mechanics, measurement on state $\Psi = \sum_i c_i \phi_i$ yields eigenstate $\phi_k$ with probability $|c_k|^2$. Formally:
\begin{equation}
\Psi \mapsto P_k \Psi = \phi_k
\end{equation}
This is instantaneous, irreversible, and non-unitary.

\subsection{Crystallisation}

We propose instead:
\begin{equation}
C[\Psi] = \Psi^{\star}
\end{equation}
where $\Psi^{\star}$ is the stable attractor of a \textbf{closure functional} $F$ over the space of admissible configurations $\mathcal{A}$:
\begin{equation}
\Psi^{\star} = \arg\min_{\Phi \in \mathcal{A}} F[\Phi]
\end{equation}
The key conceptual shift is that the definite outcome is not \textit{selected at random} but \textit{formed by constraint resolution}. What appears observationally as ``collapse'' is reinterpreted as deterministic or quasi-deterministic convergence to a constraint-compatible coherent structure.

%------------------------------------------------------------------
\section{Mathematical Framework}
%------------------------------------------------------------------

The following framework is presented as a mathematical model for state resolution, not as a derived consequence of established physical theory.

\subsection{State Space}

Let $\Psi \in \mathcal{H}_M$, the Hilbert space of configurations over a constrained manifold $M$.

\subsection{Closure Residual}

The closure residual $R[\Psi]$ measures how far the current state is from satisfying the governing constraints:
\begin{equation}
R[\Psi] = \|G[\Psi] - K\|^2 \geq 0
\end{equation}
where $G[\Psi]$ is the constraint structure induced by the state and $K$ is the admissible closure target. $R[\Psi] = 0$ indicates full coherence.

\subsection{Energy Functional}

\begin{equation}
E[\Psi] = \langle \Psi | L | \Psi \rangle
\end{equation}
where $L$ is a coupling/Laplacian/Hamiltonian-proxy operator encoding the energetic cost of the configuration.

\subsection{Coherence Metric}

\begin{equation}
Q[\Psi] = \frac{1}{n^2} \sum_{i,j} \cos(\theta_i - \theta_j)
\end{equation}
where $\theta_i = \arg(\Psi_i)$. This measures global phase alignment; $Q = 1$ indicates perfect phase coherence.

\subsection{Crystallisation Functional}

\begin{equation}
F[\Phi] = \alpha\, R[\Phi] + \beta\, E[\Phi] - \gamma\, Q[\Phi]
\end{equation}
with $\alpha, \beta, \gamma > 0$. Minimising $F$ simultaneously:
\begin{itemize}
    \item Reduces constraint mismatch ($R$)
    \item Lowers energetic cost ($E$)
    \item Rewards phase coherence ($Q$)
\end{itemize}
At this stage, the functional is introduced as a modelling assumption designed to test whether constraint-driven selection produces consistent dynamical structure.

%------------------------------------------------------------------
\section{Operator Definition}
%------------------------------------------------------------------

\begin{definition}[Crystallisation Operator]
The crystallisation operator $C$ maps an unresolved field configuration $\Psi$ on manifold $M$ to a stable coherent configuration $\Psi^{\star}$ by constrained minimisation:
\begin{equation}
C[\Psi] = \Psi^{\star} = \arg\min_{\Phi \in \mathcal{A}} F[\Phi]
\end{equation}
subject to:
\begin{equation}
\nabla_\Phi F[\Psi^{\star}] = 0, \qquad \delta^2 F[\Psi^{\star}] > 0
\end{equation}
The first condition is stationarity; the second ensures the crystallised state is a \textbf{stable} attractor, not a saddle point.
\end{definition}

%------------------------------------------------------------------
\section{Dynamical Formulation}
%------------------------------------------------------------------

Rather than instantaneous collapse, crystallisation proceeds as a gradient flow:
\begin{equation}
\frac{\partial \Psi}{\partial t} = -\nabla_\Psi F[\Psi]
\end{equation}
This means:
\begin{itemize}
    \item The state flows downhill in the constraint-energy-coherence landscape.
    \item Crystallisation is the arrival at an attractor.
    \item No discontinuity is required.
\end{itemize}

\textbf{Equilibrium:} $\nabla_\Psi F[\Psi^{\star}] = 0$

\textbf{Stability:} $\delta^2 F[\Psi^{\star}] > 0$

\subsection{Crystallisation Timescale}

Define the instability-to-resolution rate:
\begin{equation}
\Omega[\Psi] = a\,R[\Psi] + b\,Q[\Psi] + c\,|\nabla V(\Psi)|
\end{equation}
Then the crystallisation timescale is:
\begin{equation}
\tau_{\mathrm{crys}} = \frac{1}{\Omega[\Psi]}
\end{equation}
Higher mismatch or steeper landscape gradients lead to faster crystallisation. This provides a physical alternative to stochastic collapse timescales.

%------------------------------------------------------------------
\section{Spectral Formulation}
%------------------------------------------------------------------

Expand the unresolved state in eigenmodes: $\Psi = \sum_i c_i \phi_i$.

Crystallisation is not ``pick one basis vector'' but a \textbf{coherence-weighted reweighting}:
\begin{equation}
\tilde{c}_i = \frac{c_i \exp(-\mu R_i - \nu E_i + \rho Q_i)}{Z}
\end{equation}
where $Z = \sum_j c_j \exp(-\mu R_j - \nu E_j + \rho Q_j)$.

\textbf{Interpretation:}
\begin{itemize}
    \item Modes with poor closure ($R_i$ large) are suppressed.
    \item Modes with high coherence ($Q_i$ large) are amplified.
    \item The surviving structure is the coherent crystal that passes the constraint filter.
\end{itemize}

This expression should be understood as a formal reweighting model rather than a derived physical transition amplitude.

%------------------------------------------------------------------
\section{Theoretical Results}
%------------------------------------------------------------------

\begin{theorem}[Existence]
Let $F: \mathcal{A} \to \mathbb{R}$ be continuous and coercive on a non-empty, closed, bounded admissible set $\mathcal{A} \subset \mathcal{H}_M$. Then $F$ attains its minimum: there exists $\Psi^{\star} \in \mathcal{A}$ such that $F[\Psi^{\star}] \leq F[\Phi]$ for all $\Phi \in \mathcal{A}$.
\end{theorem}

\begin{proof}[Proof sketch]
$F$ is continuous and $\mathcal{A}$ is compact (closed and bounded in finite dimensions, or weakly compact in the infinite-dimensional case with coercivity). By the extreme value theorem (or direct method of calculus of variations), $F$ attains its infimum on $\mathcal{A}$. \qed
\end{proof}

\begin{theorem}[Stability]
If $\Psi^{\star}$ is a strict local minimiser of $F$ with $\nabla F[\Psi^{\star}] = 0$ and Hessian $H_F[\Psi^{\star}]$ positive definite, then $\Psi^{\star}$ is a Lyapunov-stable equilibrium of the gradient flow $\dot{\Psi} = -\nabla F[\Psi]$.
\end{theorem}

\begin{proof}[Proof sketch]
$F$ itself serves as a Lyapunov function: $\dot{F} = \langle \nabla F, \dot{\Psi}\rangle = -\|\nabla F\|^2 \leq 0$. Strict local minimality with positive-definite Hessian gives $F[\Psi] - F[\Psi^{\star}] \geq \frac{\lambda_{\min}}{2}\|\Psi - \Psi^{\star}\|^2$ in a neighbourhood, providing the required Lyapunov bound. \qed
\end{proof}

\begin{theorem}[Monotonic Descent]
Along the gradient flow $\dot{\Psi} = -\nabla F[\Psi]$:
\begin{equation}
\frac{d}{dt} F[\Psi(t)] = -\|\nabla F[\Psi(t)]\|^2 \leq 0
\end{equation}
$F$ is monotonically non-increasing.
\end{theorem}

\begin{proof}
Direct computation by chain rule. \qed
\end{proof}

\begin{theorem}[Mode Selection]
In the spectral formulation, let modes be indexed by residual $R_i$ and coherence $Q_i$. As the reweighting parameters $\mu, \rho \to \infty$ (strong constraint enforcement), the reweighted coefficients $\tilde{c}_i$ asymptotically concentrate on the mode(s) minimising $\mu R_i + \nu E_i - \rho Q_i$.
\end{theorem}

\begin{proof}[Proof sketch]
This follows from the standard Laplace/saddle-point asymptotics of exponential tilting: the mode achieving the minimum exponent dominates the partition sum $Z$ in the limit. \qed
\end{proof}

%------------------------------------------------------------------
\section{Physical Interpretation}
%------------------------------------------------------------------

The crystallisation framework provides a four-layer account of what is conventionally called ``collapse'':

\begin{table}[ht]
\centering
\caption{Four-layer account of state resolution via crystallisation.}
\label{tab:layers}
\begin{tabular}{@{}ll@{}}
\toprule
Layer & Description \\
\midrule
\textbf{Pre-state} & A manifold carries many partially compatible mode possibilities \\
\textbf{Instability} & The mixed state cannot remain globally self-consistent \\
\textbf{Crystallisation flow} & The state is driven through a closure-energy-coherence landscape \\
\textbf{Formed structure} & A coherent eigenstructure stabilises and becomes the observed state \\
\bottomrule
\end{tabular}
\end{table}

This provides a more explicit dynamical account than models in which state selection is treated as intrinsically stochastic:
\begin{enumerate}
    \item The mechanism is specified (functional minimisation).
    \item The outcome is explained (constraint-compatible attractor).
    \item The timescale is derived (not postulated).
    \item No discontinuity is required.
\end{enumerate}

%------------------------------------------------------------------
\section{Applications}
%------------------------------------------------------------------

\subsection{Quantum Measurement}

The measurement apparatus introduces boundary conditions and decoherence channels that sharpen the closure functional landscape. The system crystallises into the eigenstate(s) most compatible with the combined system--apparatus constraints. A possible interpretation is that Born probabilities may be related to the basin structure of $F$, though this connection is not established here.

\subsection{Outlook: Beyond Quantum Mechanics}

The mathematical structure of the crystallisation operator --- constrained minimisation over a functional landscape with attractor dynamics --- is not inherently quantum-mechanical. The same formalism may apply to other systems exhibiting structured state selection under constraints, including pattern formation in developmental biology and optimisation in constraint-satisfaction problems. These directions are speculative and are noted only as potential extensions; they are not developed here.

%------------------------------------------------------------------
\section{Discussion and Limitations}
%------------------------------------------------------------------

\textbf{What this framework provides:}
\begin{itemize}
    \item A precise mathematical operator providing an alternative formulation of state selection.
    \item Continuous dynamics rather than discontinuous jumps.
    \item Derived timescales rather than postulated ones.
    \item A unifying formalism with potential applicability beyond quantum mechanics.
\end{itemize}

\textbf{Current limitations:}
\begin{itemize}
    \item The choice of $\alpha, \beta, \gamma$ weights requires physical justification specific to each application domain.
    \item The connection to Born-rule statistics requires further development --- specifically, showing that basin volumes under $F$ reproduce $|c_i|^2$ weighting.
    \item Infinite-dimensional generalisation requires careful functional-analytic treatment.
\end{itemize}

\textbf{What this is not:}
\begin{itemize}
    \item A claim to replace all of quantum mechanics.
    \item A consciousness theory.
    \item A metaphysical framework.
\end{itemize}

It is strictly a mathematical and dynamical-systems reformulation of state resolution. The role of this paper is to define the operator and its mathematical properties; questions of experimental realisation and physical interpretation are addressed separately.

%------------------------------------------------------------------
\bibliographystyle{plainnat}
\bibliography{references}

\end{document}
